\chapter{طرق الغربال الأساسية وغربال سيلبرغ وعائق التكافؤ}

\section{نطاق الفصل وحالته}

\noindent\textbf{حالة الفصل:} \textenglish{AUTHORED-DRAFT / POST-AUTHORING-CORRECTIONS}. أُنشئ هذا المتن بعد
إغلاق بوابة ما قبل التأليف وصدور
\textenglish{PASS-FOR-AUTHORING = YES}. كشف تدقيق ما بعد التأليف فجوات في عرض
التصغير، وصياغة اللمّة الأساسية، وربط النتائج المقتبسة بالمراجع؛ وقد أُدخلت
التصحيحات الرياضية الأساسية، لكن التحقق المرجعي النصي المستقل ما يزال مفتوحًا.
ولا يحمل الفصل حالة \textenglish{REVIEWED} أو \textenglish{RELEASE-READY}.

يقدم الفصل الصياغة المجردة لمشكلة الغربال، ثم يبني الحد العلوي المنتهي لغربال
سيلبرغ من متراجحة مربعية ومسألة تصغير منتهية. بعد ذلك يعرّف البعد الغربالي،
ويعرض اللمّة الأساسية بوصفها مدخلًا كميًا مستقلًا، ثم يشرح عائق التكافؤ ويطبق
الحد العلوي على الأزواج المنخولة \(n,n+h\).

لا يثبت الفصل حدسية الأوليات التوأم، ولا الفجوات المحدودة، ولا حدًا سفليًا لعدد
الأزواج الأولية. كما لا يستعمل مبرهنة بومبييري--فينوغرادوف أو مبرهنة
باربان--دافنبورت--هالبرستام في بناء غربال سيلبرغ المجرد أو في التطبيق المركزي.

\subsection{نواتج التعلم}

بعد إتمام الفصل ينبغي أن يستطيع القارئ:
\begin{enumerate}
  \item صياغة بيانات الغربال بواسطة \(\mathcal A\)، و\(\mathcal P\)، و\(P(z)\)،
  و\(g(d)\)، و\(r_d\).
  \item اشتقاق متراجحة المربع التي يبدأ منها غربال سيلبرغ.
  \item فصل الصورة التربيعية الرئيسية عن حد البواقي.
  \item حل مسألة التصغير المنتهية وكتابة معاملات سيلبرغ المثلى.
  \item فهم معنى مقام سيلبرغ \(G(R,z)\) وعلاقته بحاصل الضرب المحلي \(V(z)\).
  \item التمييز بين البعد الغربالي \(\kappa\)، ومستوى البواقي \(D\)، وحد الغربلة \(z\).
  \item تفسير اللمّة الأساسية بوصفها تقريبًا أحادي الصيغة مستقلًا عن الغربال الخطي.
  \item حساب الكثافة المحلية للأزواج \(n,n+h\).
  \item صياغة عائق التكافؤ دون ادعاء استحالة مطلقة لكل طريقة غربالية.
\end{enumerate}

\section{بيانات الغربال المجردة}

\begin{definition}[بيانات الغربال]
\label{def:chapter15-sieve-data}
\resultid{ANT-DEF-15-01}
لتكن \(\mathcal A\) متتالية منتهية من أوزان غير سالبة، ولتكن \(\mathcal P\)
مجموعة من الأوليات. نضع
\[
P(z)=\prod_{\substack{p<z\\p\in\mathcal P}}p,
\qquad
S(\mathcal A,\mathcal P,z)
=
\sum_{\substack{n\in\mathcal A\\(n,P(z))=1}}1.
\]
ولكل عدد مربع حر \(d\mid P(z)\) نكتب
\[
\mathcal A_d=\{n\in\mathcal A:d\mid n\},
\qquad
|\mathcal A_d|=Xg(d)+r_d,
\]
حيث \(g\) ضربية و\(0\le g(p)<1\). ونعرّف حاصل الضرب المحلي
\[
V(z)=\prod_{\substack{p<z\\p\in\mathcal P}}(1-g(p)).
\]
\end{definition}

يمثل \(Xg(d)\) الكتلة المحلية المتوقعة للعناصر التي يقسمها \(d\)، بينما يقيس
\(r_d\) انحراف البيانات الفعلية عن هذا النموذج. لا يكفي حاصل الضرب المحلي وحده
للحصول على نتيجة كمية؛ إذ يجب أيضًا ضبط البواقي على مجال مناسب من القواسم.

\section{متراجحة سيلبرغ التربيعية}

نختار معاملات حقيقية \(\lambda_d\) مدعومة على القواسم المربعة الحرة
\(d\mid P(z)\)، مع \(\lambda_1=1\). إذا كان \((n,P(z))=1\)، فلا يقسم \(n\)
من هذه القواسم إلا \(1\)، ولذلك يساوي مجموع المعاملات \(1\).

\begin{lemma}[متراجحة المربع]
\label{lem:chapter15-square-majorant}
\resultid{ANT-LEM-15-01}
\provedhere
لدينا
\[
S(\mathcal A,\mathcal P,z)
\le
\sum_{n\in\mathcal A}
\left(
\sum_{d\mid(n,P(z))}\lambda_d
\right)^2.
\tag{15.1}
\]
\end{lemma}

\begin{proof}
إذا كان \((n,P(z))=1\)، فإن المجموع الداخلي يساوي \(\lambda_1=1\)، ومن ثم
يساوي مربعه \(1\). أما العناصر الأخرى فتضيف كميات غير سالبة، فتنتج المتراجحة
بعد الجمع على \(n\in\mathcal A\).
\end{proof}

بتوسيع المربع وتبديل المجاميع نحصل على
\[
\sum_{d_1,d_2}\lambda_{d_1}\lambda_{d_2}
|\mathcal A_{[d_1,d_2]}|
=
X\mathcal Q(\lambda)+\mathcal R(\lambda),
\tag{15.2}
\]
حيث
\[
\mathcal Q(\lambda)
=
\sum_{d_1,d_2}\lambda_{d_1}\lambda_{d_2}g([d_1,d_2]),
\]
و
\[
\mathcal R(\lambda)
=
\sum_{d_1,d_2}\lambda_{d_1}\lambda_{d_2}r_{[d_1,d_2]}.
\]

\section{حل مسألة التصغير المنتهية}

نضع للأعداد المربعة الحرة
\[
h(d)=\prod_{p\mid d}\frac{g(p)}{1-g(p)},
\qquad
f(d)=\frac1{g(d)}.
\tag{15.3}
\]
وبعكس موبيوس نكتب
\[
f(m)=\sum_{e\mid m}f_1(e).
\]
ولأن الأعداد المعنية مربعة حرة، فإن
\[
f_1(e)=\frac1{h(e)}.
\]
ولمعلمين \(R,z\ge2\) نعرّف
\[
G(R,z)=
\sum_{\substack{\ell<R\\\ell\mid P(z)}}h(\ell).
\tag{15.4}
\]

\begin{theorem}[الحد العلوي المنتهي لغربال سيلبرغ]
\label{thm:chapter15-finite-selberg-upper}
\resultid{ANT-THM-15-01}
\provedhere
إذا كانت معاملات الغربال مدعومة على \(d<R\)، فإن هناك اختيارًا للمعاملات
يحقق
\[
S(\mathcal A,\mathcal P,z)
\le
\frac{X}{G(R,z)}
+
\sum_{\substack{d<R^2\\d\mid P(z)}}3^{\omega(d)}|r_d|.
\tag{15.5}
\]
\end{theorem}

\begin{proof}
بما أن \(g\) ضربية وعلى الأعداد المربعة الحرة، فإن
\[
g([d_1,d_2])
=
\frac{g(d_1)g(d_2)}{g((d_1,d_2))}.
\]
ومن ثم
\[
\mathcal Q(\lambda)
=
\sum_{d_1,d_2<R}
\lambda_{d_1}\lambda_{d_2}
\frac{f((d_1,d_2))}{f(d_1)f(d_2)}.
\]
نعرف المتغيرات المحولة
\[
u_e=
\sum_{\substack{e\mid d<R\\d\mid P(z)}}
\frac{\lambda_d}{f(d)}.
\]
وباستخدام \(f(m)=\sum_{e\mid m}f_1(e)\) نحصل على التقطير
\[
\mathcal Q(\lambda)
=
\sum_{\substack{e<R\\e\mid P(z)}}f_1(e)u_e^2
=
\sum_{\substack{e<R\\e\mid P(z)}}\frac{u_e^2}{h(e)}.
\tag{15.5a}
\]
والتحويل مثلثي، ويعكسه موبيوس على الصورة
\[
\lambda_d
=
f(d)
\sum_{\substack{d\mid e<R\\e\mid P(z)}}
\mu(e/d)u_e.
\tag{15.5b}
\]
وعند \(d=1\) يصبح القيد \(\lambda_1=1\)
\[
\sum_{\substack{e<R\\e\mid P(z)}}\mu(e)u_e=1.
\]
تعطي متراجحة كوشي--شفارتس
\[
1
\le
\left(\sum_e\frac{u_e^2}{h(e)}\right)
\left(\sum_e h(e)\right)
=
\mathcal Q(\lambda)G(R,z),
\]
لأن \(\mu(e)^2=1\) على القواسم المربعة الحرة. إذن
\[
\mathcal Q(\lambda)\ge\frac1{G(R,z)}.
\]
ويتحقق التساوي عند
\[
u_e=\frac{\mu(e)h(e)}{G(R,z)}.
\]
وبالتعويض في (15.5b) نجد معاملات سيلبرغ المثلى
\[
\boxed{
\lambda_d
=
\frac{\mu(d)}{g(d)G(R,z)}
\sum_{\substack{d\mid e<R\\e\mid P(z)}}h(e)
}.
\tag{15.5c}
\]
أما حد البواقي، فلكل عدد مربع حر \(d\)، يساوي عدد الأزواج المرتبة
\((d_1,d_2)\) التي تحقق \([d_1,d_2]=d\) بالضبط \(3^{\omega(d)}\): لكل أولي
يقسم \(d\) ثلاثة اختيارات، أن يظهر في \(d_1\) فقط، أو في \(d_2\) فقط، أو
فيهما معًا. كما أن \(d_1,d_2<R\) يقتضيان \([d_1,d_2]<R^2\). ينتج الحد
المعلن من (15.1)--(15.2).
\end{proof}

\section{البعد الغربالي ومقام سيلبرغ}

نفرض شرط الانتظام التالي: يوجد ثابت \(A_1\ge0\) بحيث لكل \(2\le w<z\)،
\[
\left|
\sum_{w\le p<z}g(p)\log p
-
\kappa\log\frac zw
\right|
\le A_1.
\tag{15.6}
\]
وعندئذ يكون
\[
V(z)\asymp_{\kappa,A_1}\frac1{(\log z)^\kappa}.
\]
يجب عدم الخلط بين \(\kappa\) وبين مستوى البواقي \(D\)، ولا بينهما وبين حد
الغربلة \(z\).

\begin{theorem}[تقدير مقام سيلبرغ في البعد \(\kappa\)]
\label{thm:chapter15-selberg-denominator}
\resultid{ANT-THM-15-02}
\citedresult{\cite[p.~21]{HeathBrown2002LecturesSieves}; \cite[(5.3.1)]{HalberstamRichert1974}}
إذا تحقق شرط الانتظام (15.6)، فإن
\[
G(z,z)
=
\frac{e^{\gamma\kappa}}{\Gamma(\kappa+1)}
V(z)^{-1}
\left(1+O_{\kappa,A_1}\!\left(\frac1{\log z}\right)\right).
\tag{15.7}
\]
\end{theorem}

الثابت في (15.7) متوافق مع حالة \(\kappa=1\)، إذ يعيد ثابت ميرتنز
\(e^\gamma\). أما تحديد رقم المبرهنة والصفحات في المرجع المختار فيبقى بندًا
مفتوحًا في التدقيق المرجعي النصي، ولذلك لا يرفع هذا الفصل بعد إلى
\textenglish{REFERENCE-AUDIT = PASS}.

\section{اللمّة الأساسية وحدودها}

إذا كانت البواقي مضبوطة حتى مستوى \(D\)، نضع
\[
s=\frac{\log D}{\log z}.
\tag{15.8}
\]
اللمّة الأساسية ليست هي مبرهنة الغربال الخطي ذات الدالتين
\(F_\kappa,f_\kappa\). إنها تقريب للكمية المنخولة بالحد المحلي \(XV(z)\) عندما
يكبر معامل الغربلة \(s\)، مع خطأ يضم جزءًا بنيويًا يتناقص مع \(s\) وجزءًا
قادما من البواقي.

\begin{theorem}[صيغة عامة للّمة الأساسية]
\label{thm:chapter15-fundamental-lemma}
\resultid{ANT-THM-15-03}
\citedresult{\cite[Chapter~4, pp.~29--42]{DiamondHalberstamGalway2008}}
تحت شرط الانتظام (15.6)، توجد عتبة
\(s_0=s_0(\kappa,A_1)\) ودالة غير سالبة
\(\rho_{\kappa,A_1}(s)\to0\) عندما \(s\to\infty\)، بحيث إذا
\(s\ge s_0\) وكانت البواقي مضبوطة حتى \(D\)، فإن نسخة اللمّة الأساسية المختارة
تعطي
\[
S(\mathcal A,\mathcal P,z)
=
XV(z)
\left(1+O_{\kappa,A_1}(\rho_{\kappa,A_1}(s))\right)
+
O\!\left(
\sum_{\substack{d<D\\d\mid P(z)}}
 c_\kappa^{\omega(d)}|r_d|
\right),
\tag{15.9}
\]
حيث \(c_\kappa\ge1\) ثابت تحدده صياغة اللمّة المستعملة.
\end{theorem}

لا نثبت هنا معدلًا عدديًا معينًا لـ\(\rho_{\kappa,A_1}\)، لأن ذلك يتطلب اختيار
نسخة مرجعية بعينها وتثبيت تطبيعها ومجالها. كما لا نستعمل الرمزين
\(F_\kappa,f_\kappa\) هنا؛ فهما ينتميان إلى صيغ الغربال الخطي أو
\textenglish{beta-sieve} المتخصصة، وهي مؤجلة إلى فصل لاحق.

\section{تطبيق: الأزواج المنخولة}

لتكن \(h\ne0\) ثابتة، ولتكن
\[
S_h(x,z)
=
\#\{1\le n\le x:(n(n+h),P(z))=1\}.
\tag{15.10}
\]
ولكل أولي \(p\)، نعرّف \(\rho_h(p)\) بعدد حلول التطابق
\[
n(n+h)\equiv0\pmod p.
\]

\begin{proposition}[الكثافة المحلية للأزواج]
\label{prop:chapter15-pair-local-density}
\resultid{ANT-PROP-15-01}
\provedhere
لدينا
\[
\rho_h(p)=
\begin{cases}
2,&p\nmid h,\\
1,&p\mid h.
\end{cases}
\tag{15.11}
\]
ومن ثم \(g_h(p)=\rho_h(p)/p\). وإذا كان \(h\) فرديًا و\(z>2\)، فإن
\(S_h(x,z)=0\).
\end{proposition}

\begin{proof}
الحلول هي الفئتان \(n\equiv0\pmod p\) و\(n\equiv-h\pmod p\). تتطابقان إذا
وفقط إذا \(p\mid h\). وإذا كان \(h\) فرديًا، فإن أحد العددين \(n,n+h\) زوجي
دائمًا، ولذلك يقسم \(2\) حاصل الضرب عندما تدخل الأولية \(2\) في \(P(z)\).
\end{proof}

\begin{proposition}[البعد الغربالي للأزواج]
\label{prop:chapter15-pair-dimension}
\resultid{ANT-PROP-15-02}
\provedhere
إذا كان \(h\ne0\) زوجيًا، فإن بيانات الأزواج ذات بعد غربالي \(2\)، ولدينا
\[
V_h(z)
=
\prod_{p<z}\left(1-\frac{\rho_h(p)}p\right)
\asymp_h \frac1{(\log z)^2}.
\tag{15.12}
\]
\end{proposition}

\begin{proof}
باستثناء الأوليات القاسمة لـ\(h\)، يكون العامل المحلي \(1-2/p\). نقارن
حاصل الضرب هذا بمربع حاصل ضرب ميرتنز
\(\prod_{p<z}(1-1/p)^2\). نسبة العاملين تساوي \(1+O(p^{-2})\)، ولذلك يتقارب
حاصل ضرب النسب. أما الأوليات القاسمة لـ\(h\) فتعدل الثابت فقط.
\end{proof}

للقواسم المربعة الحرة \(d\)، تعطي مبرهنة الباقي الصيني
\[
|\mathcal A_d|
=
\frac{x\rho_h(d)}d+O(\rho_h(d)),
\qquad
\rho_h(d)=\prod_{p\mid d}\rho_h(p),
\tag{15.13}
\]
ومن ثم يمكن أخذ \(X=x\)، و\(g_h(d)=\rho_h(d)/d\)، و
\(|r_d|\le \rho_h(d)\le2^{\omega(d)}\).

\begin{theorem}[حد الأزواج المنخولة]
\label{thm:chapter15-sifted-pairs}
\resultid{ANT-THM-15-04}
\provedhere
لكل ثابت زوجي \(h\ne0\)، وبانتظام عندما
\[
x\ge3,
\qquad
3\le z\le x^{1/4},
\]
لدينا
\[
S_h(x,z)
\ll_h
\frac{x}{(\log z)^2}.
\tag{15.14}
\]
\end{theorem}

\begin{proof}
نطبق المبرهنة \ref{thm:chapter15-finite-selberg-upper} مع \(R=z\). من
المبرهنة \ref{thm:chapter15-selberg-denominator} والقضية
\ref{prop:chapter15-pair-dimension} نحصل على
\[
\frac{x}{G_h(z,z)}\ll_h \frac{x}{(\log z)^2}.
\]
أما حد البواقي فيحقق
\[
\sum_{d<z^2}3^{\omega(d)}|r_d|
\le
\sum_{d<z^2}6^{\omega(d)}
\ll z^2(\log z)^5.
\]
والشرط الكمي الكافي لامتصاص هذا الحد هو
\[
z^2(\log z)^7\ll x.
\]
اختيار \(z\le x^{1/4}\) مجال تعليمي آمن ومتحفظ، وليس الأس الحرج الأمثل. ولـ
\(x\) الأكبر من ثابت مطلق مناسب يعطي هذا الاختيار
\[
z^2(\log z)^5
\ll
\frac{x}{(\log z)^2}.
\]
أما المجال المحدود المتبقي، فنستخدم الحد التافه \(S_h(x,z)\le x\). وبما أن
\(z\ge3\)، فإن \((\log z)^2\ge(\log3)^2>0\)، فيمتص عدد محدود من القيم في
الثابت الضمني المعتمد على \(h\). ينتج (15.14).
\end{proof}

\section{عائق التكافؤ}

\begin{remark}[عائق التكافؤ]
\label{rem:chapter15-parity-barrier}
\resultid{ANT-DIAG-15-01}
\citedresult{\cite[Chapter~6]{IwaniecKowalski2004}; \cite{FriedlanderIwaniec1998AsymptoticSieve}}
تعتمد المناخل التقليدية أساسًا على بيانات القواسم المحلية من نوع
\(|\mathcal A_d|=Xg(d)+r_d\). هذه البيانات لا تميز وحدها، في العموم، بين
العناصر ذات عدد زوجي أو فردي من العوامل الأولية. لذلك تستطيع المناخل إنتاج
حدود عليا قوية ونتائج عن الأعداد شبه الأولية، لكنها تصطدم بحاجز بنيوي عند
محاولة استخراج حد سفلي دقيق لوجود أوليات ضمن نمط إضافي.
\end{remark}

لا تعني هذه الملاحظة أن كل طريقة يدخل فيها غربال عاجزة مطلقًا. فقد تتجاوز
بعض التطبيقات صورًا من الحاجز بإضافة بنية لا توفرها بيانات القواسم المحلية
وحدها، مثل الهويات الالتفافية، والأشكال الثنائية، ومعلومات توزيع أعمق، أو
مداخل طيفية. وفي تطبيق الأزواج المنخولة، يفسر العائق لماذا لا يتحول الحد
العلوي (15.14) إلى برهان على وجود أزواج أولية.

\section{ملاحظات ختامية وتمارين}

يلخص هذا الفصل الطبقات التي ينبغي عدم خلطها:
\begin{enumerate}
  \item متراجحة المربع وحل التصغير المنتهي: مثبتان هنا.
  \item التحليل التقاربي لمقام سيلبرغ: مدخل مقتبس، وموضعه النصي الدقيق قيد التدقيق.
  \item اللمّة الأساسية: تقريب أحادي الصيغة، لا دالتا الغربال الخطي.
  \item الغربال الخطي والـ\textenglish{beta-sieve}: أدوات متخصصة مؤجلة.
  \item عائق التكافؤ: تشخيص لحدود المعلومات المحلية، لا استحالة مطلقة.
\end{enumerate}

\begin{exercise}
تحقق مباشرة من أن عدد الأزواج المرتبة \((d_1,d_2)\) من القواسم المربعة الحرة
التي تحقق \([d_1,d_2]=d\) يساوي \(3^{\omega(d)}\).
\end{exercise}

\begin{exercise}
أثبت أن
\[
\prod_{\substack{p<z\\p\nmid h}}
\frac{1-2/p}{(1-1/p)^2}
\]
يتقارب إلى ثابت موجب عندما \(h\) ثابت زوجي غير صفري.
\end{exercise}

\begin{exercise}
فسر لماذا لا يكفي الحد العلوي
\(S_h(x,z)\ll_h x/(\log z)^2\) لإثبات وجود عدد لا نهائي من الأزواج الأولية
\(p,p+h\).
\end{exercise}

\begin{exercise}
قارن بين الأدوار المختلفة للمعلمات \(z\)، و\(R\)، و\(D\)، و\(\kappa\) في
الفصل، واذكر موضع ظهور كل واحدة منها.
\end{exercise}
